\documentclass[12pt,a4paper]{article}
\usepackage[utf8]{inputenc}
\usepackage[T1]{fontenc}
\usepackage{amsmath,amssymb,amsfonts}
\usepackage{geometry}
\usepackage{booktabs}
\usepackage{hyperref}
\geometry{margin=2.5cm}

\title{Ontology of Causally Self-Consistent Regions}
\author{Kevin Mu\~noz}
\date{}

\begin{document}
\maketitle

\begin{abstract}
We develop an ontological articulation compatible with a horizon-centered approach to cosmology and gravitational physics. Rather than introducing a new physical theory or proposing additional dynamical laws, the present work explores the minimal ontological commitments required by frameworks in which causal accessibility and horizon structure play a structurally primary role.

Within this perspective, we introduce the notion of a causally self-consistent region (CSCR): a physical domain characterized by causal accessibility and bounded by a horizon structure that constrains the relevant degrees of freedom available to internal observers.

The concept is intended as an ontological clarification rather than as a complete model of reality. It does not assume a specific microscopic substrate, global spacetime completion, or unique cosmological realization. Instead, it provides a minimal conceptual structure capable of organizing horizon-centered descriptions in a manner consistent with gravitational thermodynamics and causal frameworks.

The interior of black-hole geometries is discussed only as an illustrative realization and not as a literal identification.
\end{abstract}

\section{Introduction}

Developments in gravitational thermodynamics and horizon physics increasingly suggest that causal boundaries possess structural significance beyond their traditional geometric interpretation. Black hole thermodynamics, cosmological horizons, and causal-structure-based approaches consistently indicate that entropy, information accessibility, and observer-dependent descriptions may be constrained by horizons rather than determined solely by volumetric properties.

Motivated by these observations, horizon-centered frameworks have been proposed in which causal structure is treated as an organizing principle for cosmological interpretation.

The present work does not extend or modify such frameworks. Instead, it examines what minimal ontological commitments are required if causal accessibility and horizon structure are regarded as physically significant.

The purpose is not to provide a fundamental ontology of reality, but rather to articulate a domain concept capable of supporting horizon-centered descriptions while remaining compatible with multiple effective realizations.

\section{Status and Scope of Claims}

The framework developed here is ontological and conceptual in character.

It does not introduce:

\begin{itemize}
  \item new dynamical laws
  \item new cosmological models
  \item microscopic descriptions
  \item modifications of General Relativity
  \item observational predictions
\end{itemize}

The notion of a causally self-consistent region should therefore be interpreted as a minimal consistency structure rather than as an independent physical theory.

No assumption is made concerning:

\begin{itemize}
  \item the totality of spacetime
  \item the existence of external domains
  \item discreteness of underlying degrees of freedom
  \item quantum gravitational mechanisms
  \item fundamental reduction of geometry
\end{itemize}

Statements concerning emergence and ontology should therefore be interpreted at the level of structural description rather than metaphysical assertion.

\section{Causally Self-Consistent Regions}

\subsection{Core Definition}

A causally self-consistent region (CSCR) is defined as the minimal physical domain required for a consistent description in terms of causal accessibility and horizon structure.

A CSCR is characterized by:

\begin{itemize}
  \item causal accessibility among internal degrees of freedom
  \item a horizon-defined boundary
  \item operationally meaningful observables
  \item internal consistency between accessible information and effective dynamics
\end{itemize}

The concept does not imply that such regions constitute the entirety of physical reality.

Rather, they represent the minimal domains within which horizon-centered descriptions become operationally meaningful.

No assumption is made regarding embedding within larger structures.

In the cosmological effective realization, the CSCR boundary is located at the apparent horizon radius $R_A = 1/H$, and perturbations of the boundary correspond to $\psi = \delta R_A / R_A$.

\subsection{Ontological Principles}

\paragraph{Causal Primacy:} Causal relations are treated as structurally prior to geometric description. Spatial geometry is regarded as an effective representation of causal accessibility rather than an ontologically fundamental entity. No particular mechanism for such emergence is assumed.

\paragraph{Operational Causal Closure:} Observable physical influence does not cross the horizon within the domain of description. Statements concerning external regions are regarded as undefined rather than false. This principle should be interpreted operationally and does not imply absolute metaphysical isolation.

\paragraph{Horizon-Defined Domain:} The horizon determines the physically relevant domain available to internal observers. The horizon is therefore treated as a defining structure of the description rather than merely as a geometric surface. No claim is made regarding the microscopic constitution of horizon degrees of freedom.

\paragraph{Boundary-Associated Entropy:} Entropy relevant to the domain is assumed to be primarily associated with horizon structure. This principle follows established insights from gravitational thermodynamics and should be understood as a structural guideline rather than a derivation. No quantitative entropy model is introduced.

\paragraph{Consistency Constraint:} Any admissible effective description of a CSCR should remain compatible with the causal and entropic properties of its defining boundary. Bulk descriptions are therefore interpreted as constrained by boundary structure. No unique dynamical realization is implied.

In the flat cosmological realization, this constraint is exemplified by the relation $\delta\rho/\rho = -2\psi$: density perturbations in the bulk are determined by fractional fluctuations of the horizon boundary, making explicit how boundary structure constrains bulk descriptions.

\section{Emergence of Space and Time}

Within horizon-centered frameworks, geometry and temporal ordering need not be regarded as fundamental.

Spatial relations may instead be interpreted as effective descriptions of causal accessibility among degrees of freedom.

Similarly, temporal ordering may arise from changes in available information and evolving causal relations.

The present work does not attempt to derive these structures.

Rather, it identifies conceptual compatibility with approaches in which spacetime is emergent.

\section{Realizations and Illustrative Examples}

The CSCR concept does not prescribe a unique spacetime realization.

Different physical settings may instantiate compatible structures, including:

\begin{itemize}
  \item cosmological horizon systems
  \item horizon-dominated phases
  \item black-hole interiors
  \item causal-structure-based constructions
\end{itemize}

These examples should be understood as illustrative realizations rather than defining cases.

Compatibility does not imply equivalence.

\section{Illustrative Case: Black-Hole Interior Geometry}

The interior region of a black-hole spacetime provides a particularly useful conceptual example.

Such systems possess:

\begin{itemize}
  \item causal boundaries
  \item observer-dependent descriptions
  \item horizon-associated entropy
  \item constrained accessibility
\end{itemize}

These features resemble those required of a causally self-consistent region.

This comparison is intended solely as a conceptual analogy.

No identification between cosmological domains and black-hole interiors is proposed.

No conclusions concerning cosmological origin follow from this example.

\section{Open Directions}

Several questions remain open:

\begin{itemize}
  \item microscopic origin of horizon degrees of freedom
  \item quantitative boundary-to-bulk relations
  \item emergence of effective geometry
  \item relation to quantum gravity
  \item possible effective realizations
\end{itemize}

These open questions define future research directions rather than deficiencies.

\section{Conclusion}

We have introduced the notion of causally self-consistent regions as a minimal ontological articulation compatible with horizon-centered approaches.

The purpose of the framework is not to provide a complete ontology or replace established physical descriptions.

Instead, it offers a minimal conceptual structure capable of supporting interpretations organized around causal accessibility, horizon structure, and boundary-associated entropy.

Future developments will determine whether these ideas admit quantitative realization within broader theoretical frameworks.

\end{document}
